\chapter{Xem Đáp Án Trước Khi Hiểu}
\chapterintrobi{Muốn ra kết quả nhanh hơn muốn đi qua con đường hiểu bài}{Look at the answer first}{Một số bạn yêu cảm giác biết kết quả hơn yêu quá trình hiểu vì sao ra kết quả đó.}

\begin{lucbatbox}{Lục bát: Biết rồi nhưng chưa hiểu}
Bài kia em mới đọc qua\
Đáp án em đã tìm ra tức thì\
Nhìn xong thấy cũng an yên\
Tưởng mình hiểu hết, hóa phiền về sau\
Hỏi vì sao lại ra màu\
Em đành im lặng, cúi đầu với trang
\end{lucbatbox}

\begin{tutuyetbox}{Thất ngôn tứ tuyệt: Kết quả đi trước}
Đáp án hiện lên rất nhẹ nhàng\
Em mừng như thể đã vỡ vàng\
Nhưng phần bản chất chưa hề mở cửa\
Nên mai gặp lại vẫn ngỡ ngàng
\end{tutuyetbox}

\begin{ghichu}
Xem đáp án trước khi hiểu giúp bớt sốt ruột trước mắt, nhưng thường làm mình vay sự rõ ràng của người khác mà chưa tự có.
\end{ghichu}

\begin{englishnote}
\textbf{answer key}: đáp án.\par
\textbf{understand the process}: hiểu cách làm.\par
\textbf{I saw the answer before I understood it.}: Em xem đáp án trước khi thật sự hiểu.
\end{englishnote}